% =============================================================================
% CompetitiveExam/CS301/05-queues-questions.tex
% Competitive Exam Practice 5 — GATE (Computer Science) Pattern
% Queue ADT and linear-array false-full trap, circular queue with mod
% wraparound and full/empty tests, double-ended queue and restricted
% variants, priority queue using arrays (unordered vs ordered bills),
% plus Unit-I review (Labs P5, P6; Class Test 1).
%
% Student-facing — NO answers. Answer key: 05-queues-answers.tex.
%
% Provenance (availability-first, Tier 1 = local PYQ book
% CompetitiveExam/Books/GATE-PYQs/filter1_volume2.pdf):
%   - Q1  [GATE CSE 1996 Q1.12, verified — GATEOverflow V2 p.240] —
%     FIFO/LIFO-by-stacks-or-queues + list/queue-implementation
%     statements. Stem fully recoverable via pdftotext; option numerals
%     image-dropped locally, recovered verbatim from the BYJUS GATE
%     question bank ((a) ii and iii / (b) i and ii / (c) iii and iv /
%     (d) ii and iv; "Answer (a)"); year confirmed on ExamSIDE's
%     GATE-1996 listing. Only (ii) and (iii) are true, by direct
%     evaluation (stacks are LIFO, queues are FIFO; circular reuse
%     beats linear drift).
%   - Q2  [GATE CSE 2012 Q35, verified — GATEOverflow V2 p.241] —
%     circular-queue full/empty conditions with REAR/FRONT, capacity
%     n-1 in n slots. Stem + all four options confirmed on the
%     GATEOverflow question page (32k views), Testbook's GATE-2012
%     solved page, and ExamSIDE's GATE-2012 listing; answer (A)
%     confirmed on Testbook ("option 1") and the GATEOverflow
%     accepted reasoning (initial REAR = FRONT = 0 kills (B)/(C)/(D)).
%   - Q3  [GATE CSE 2016 Set 1 Q10, verified — GATEOverflow V2 p.242]
%     — array queue with efficient ENQUEUE and DEQUEUE: both O(1).
%     Stem + all four options confirmed on appliedroots' GATE-2016-Set-1
%     listing and the SamagraCS solved-PYQ page, and listed on
%     ExamSIDE's stacks-and-queues topic page; answer (A) per the
%     official GATE 2016 key and deterministic from the deck's circular
%     implementation (one array access + one mod per op, O(1) worst).
%   - Q4  [GATE CSE 1997 Q4.7, verified — GATEOverflow V2 p.239] —
%     priority queue used as a stack (PUSH = INSERT with key K,
%     POP = DELETEMIN): keys in strictly decreasing order. Stem +
%     all four options fully recoverable via pdftotext (Tier-1 text);
%     listed on ExamSIDE's stacks-and-queues page under GATE CSE 1997;
%     answer (D) derived deterministically (DELETEMIN returns the
%     smallest key, so LIFO needs each newer key smaller than all
%     older ones) and consistent with the documented monotonic-key
%     principle (StackOverflow: stack-from-priority-queue uses
%     monotonic priorities; min-extraction forces decreasing).
%   - Q5  [GATE IT 2007 Q30, verified — GATEOverflow V2 pp.243-244]
%     — recursive f (delete, recurse, insert) reverses the queue.
%     Stem + all four options fully recoverable via pdftotext;
%     confirmed on the GATEOverflow question page (23.8k views, best
%     answer B with 56 votes), the GFG GATE-IT-2007 quiz, and
%     ExamSIDE's stacks-and-queues listing. Answer (B).
%   - Q6  [Original, GATE-pattern] — circular-queue count-version
%     numeric trace (capacity 6, front = 4, rear = 2, count = 5;
%     dequeue then enqueue). Fresh instance: the deck's worked trace
%     uses capacity 4 (front 2/rear 2/count 1 walk) and its exit ticket
%     uses capacity 5 (front 3/rear 1/count 4). Verified by execution
%     2026-09-17: final front = 5, rear = 3, count = 5, not full.
%   - Q7  [Original, GATE-pattern] — restricted-deque naming: insert
%     at both ends, delete front-only => output-restricted; forbidden
%     op is dequeueRear. Fresh instance (the deck's worked example
%     uses enqueueRear(Q)/enqueueRear(R)/enqueueFront(P); no prior
%     practice set tests the naming). Cross-checked against the
%     standard definitions (input-restricted = single-end insert;
%     output-restricted = single-end delete) and the BU Lecture-5
%     restricted-deque operation sets. NOTE (open issue, flagged in
%     the gate report, not fixed here): the deck's exam-note box in
%     Act 3 labels "insert at rear, delete at both ends" as
%     output-restricted, which contradicts both the standard
%     definitions and the deck's own bullet definitions two paragraphs
%     earlier (insert-one-end-only = input-restricted). Q7 follows the
%     standard/BU naming.
%   - Q8  [Original, GATE-pattern] — unordered vs ordered array
%     priority-queue workload choice (10,000 inserts then 5
%     deleteMins => unordered wins). Fresh instance: the deck's
%     check-your-prediction uses 1,000/3 and 3/1,000 splits. Verified
%     by execution 2026-09-17 (unordered ~60k vs ordered ~1e8
%     worst-case elementary ops).
%
% Considered and dropped: GATE CSE 2022 Q52 (two-queue reversal NAT,
% V2 p.242-243) — figure-dependent (Initial/Final-state diagrams are
% embedded images, dropped from the text layer) with no second web
% source confirming exact queue contents; GATE CSE 2016 Set 1 Q41
% (queue+stack sort NAT, V2 p.242) — a sorting-network question, out
% of Week 5 array-queue scope; GATE DS&AI 2024 Q22 (deque op trace,
% V2 p.243) — operation lines truncated in the text layer
% ("a removeFirst()" fragments) with no second source confirming the
% full stem, so not shippable as verified (its style is mirrored by
% original Q6/Q7 instead); GATE CSE 2017 Set 2 Q13 and GATE CSE 2018
% Q3 (linked-list queues) — Week 7 scope, not arrays; GATE CSE 2006
% Q49 and GATE CSE 2013 Q44 (amortized two-stack/MultiDequeue
% analysis) — beyond Week 5's O(1)-per-op array scope.
%
% No render+OCR pass was possible (no Ghostscript/Tesseract in this
% environment); image-dropped values above were recovered via Tier-2 web
% search + CoVe instead, per the skill's sourcing ladder. No forked
% claim-verifier subagent exists in this environment, so CoVe was executed
% inline (each claim checked against 2+ independent sources, verdicts in
% the session report).
%
% Compile with: cd CompetitiveExam/CS301 &&
%   TEXINPUTS=../../Preambles;$TEXINPUTS xelatex -interaction=nonstopmode 05-queues-questions.tex
%   (Windows/MiKTeX uses ; not : as the separator)
% =============================================================================

\documentclass[11pt]{article}
\usepackage[margin=1in]{geometry}
% =============================================================================
% Preambles/header.tex — shared LaTeX/Beamer preamble for this project
%
% Use from a Beamer lecture by prepending:
%     \input{header}
% and compiling with TEXINPUTS=../Preambles:$TEXINPUTS (the /compile-latex
% skill does this automatically).
%
% PALETTE CONTRACT. Color names defined here MUST match the names in
% Quarto/theme-template.scss so Beamer and Quarto renderings use the same
% palette. The scripts/check-palette-sync.sh script enforces this.
%
% To customize: edit the HEX values below AND the matching $variable in
% Quarto/theme-template.scss, then run scripts/check-palette-sync.sh.
% =============================================================================

% -----------------------------------------------------------------------------
% Engine detection + encoding
%
% Our /compile-latex skill uses XeLaTeX, where inputenc is unnecessary and
% can produce encoding warnings. Only load inputenc under pdfTeX.
% -----------------------------------------------------------------------------
\usepackage{iftex}
\ifPDFTeX
  \usepackage[utf8]{inputenc}
\fi

\usepackage{amsmath, amssymb, amsthm}
\usepackage{booktabs}
\usepackage{graphicx}

% -----------------------------------------------------------------------------
% PALETTE — mirrors Quarto/theme-template.scss variable names.
% Load xcolor and define colors BEFORE anything that references them
% (notably hyperref, hypersetup, and the Beamer color assignments).
% -----------------------------------------------------------------------------
\usepackage{xcolor}

\definecolor{primary-blue}{HTML}{012169}      % headings, accents
\definecolor{primary-gold}{HTML}{B9975B}      % emphasis, borders
\definecolor{highlight-yellow}{HTML}{F2A900}  % markers, alerts
\definecolor{light-bg}{HTML}{E8EDF5}          % title / section backgrounds
\definecolor{jet}{HTML}{1A1A1A}               % body text

% Semantic colors (used in diagrams and annotations)
\definecolor{positive}{HTML}{15803D}          % good / observed / identified
\definecolor{negative}{HTML}{B91C1C}          % bad / problematic / bias
\definecolor{neutral}{HTML}{525252}           % reference / context

% Highlight accents (match .hi-* classes in the SCSS)
\definecolor{hi-slate}{HTML}{314F4F}
\definecolor{hi-green}{HTML}{2E8B57}
\definecolor{hi-red}{HTML}{C41E3A}

% Semantic aliases so skill templates and snippets can write
% \textcolor{accent}{...} without hard-coding "primary-blue".
\colorlet{accent}{primary-blue}
\colorlet{accent2}{primary-gold}

% -----------------------------------------------------------------------------
% Hyperref — loaded AFTER xcolor + palette so link colors resolve.
% -----------------------------------------------------------------------------
\usepackage{hyperref}
\hypersetup{colorlinks=true, linkcolor=primary-blue, urlcolor=primary-blue,
            citecolor=primary-blue, pdfborder={0 0 0}}

% -----------------------------------------------------------------------------
% Course code — set once per deck (after \input{header}, before \begin{document})
% via \coursecode{CS401}. Shown in the footer on every slide so a deck's course
% is identifiable without opening the title page. Empty by default.
% -----------------------------------------------------------------------------
\makeatletter
\newcommand{\coursecode}[1]{\gdef\@coursecodetext{#1}}
\gdef\@coursecodetext{}
\makeatother

% -----------------------------------------------------------------------------
% Beamer theme assignments (only apply under Beamer).
%
% We detect Beamer via \@ifclassloaded{beamer}, which is the documented,
% non-internal way. This must be wrapped in \makeatletter/\makeatother
% because \@ifclassloaded contains an @-catcode character.
% -----------------------------------------------------------------------------
\makeatletter
\@ifclassloaded{beamer}{%
  \usefonttheme{professionalfonts}
  \setbeamercolor{structure}{fg=primary-blue}
  \setbeamercolor{frametitle}{fg=primary-blue}
  \setbeamercolor{title}{fg=primary-blue}
  \setbeamercolor{subtitle}{fg=primary-gold}
  \setbeamercolor{author}{fg=primary-blue}
  \setbeamercolor{institute}{fg=neutral}
  \setbeamercolor{date}{fg=primary-gold}
  \setbeamercolor{itemize item}{fg=primary-blue}
  \setbeamercolor{itemize subitem}{fg=primary-gold}
  \setbeamercolor{alerted text}{fg=primary-gold}
  \setbeamercolor{block title}{fg=white, bg=primary-blue}
  \setbeamercolor{block body}{bg=light-bg}
  % Semantic block variants: block=definition/concept (blue), exampleblock=
  % worked example (green/positive), alertblock=key takeaway or caution (gold).
  % Keep to <=2 colored boxes per slide across ALL three variants (INV-7).
  \setbeamercolor{block title example}{fg=white, bg=positive}
  \setbeamercolor{block body example}{bg=light-bg}
  \setbeamercolor{block title alerted}{fg=white, bg=primary-gold}
  \setbeamercolor{block body alerted}{bg=light-bg}

  % Minimal footer: course code (left) + page X / N (right), no navigation clutter
  \setbeamertemplate{navigation symbols}{}
  \setbeamertemplate{footline}{%
    \color{neutral}\footnotesize\hspace{0.7em}\@coursecodetext\hfill
    \insertframenumber/\inserttotalframenumber\hspace{0.7em}\vspace{0.3em}}
}{}
\makeatother

% -----------------------------------------------------------------------------
% TikZ setup — libraries the snippet gallery uses
% -----------------------------------------------------------------------------
\usepackage{tikz}
\usetikzlibrary{arrows.meta, positioning, calc, decorations.pathreplacing,
                fit, shapes.geometric, backgrounds}

% Shared TikZ styles — snippets and hand-written diagrams can pick these up
% via `\begin{tikzpicture}[every node/.style={...}]` or by naming specific
% styles (e.g., `\node[dag-node] (X) at (0,0) {X};`).
\tikzset{
  % Node styles for causal DAGs and similar
  dag-node/.style = {
    draw=primary-blue, fill=light-bg, rounded corners=2pt,
    minimum width=1.2cm, minimum height=0.9cm, align=center, font=\small
  },
  decision-node/.style = {
    draw=primary-gold, fill=white, diamond, aspect=1.8,
    minimum width=1.8cm, minimum height=1.2cm, align=center, font=\small,
    inner sep=1pt
  },
  % Edge styles — observed vs counterfactual
  observed-edge/.style = {-{Latex[length=2.5mm]}, thick, draw=jet},
  counterfactual-edge/.style = {-{Latex[length=2.5mm]}, thick, draw=neutral, dashed},
  confound-edge/.style = {-{Latex[length=2.5mm]}, thick, draw=negative, dashed},
  % Data-point styles for plots
  observed-dot/.style = {fill=primary-blue, draw=primary-blue, circle, inner sep=1.2pt},
  counterfactual-dot/.style = {fill=white, draw=neutral, circle, inner sep=1.2pt}
}

% -----------------------------------------------------------------------------
% Convenience macros
% -----------------------------------------------------------------------------
\newcommand{\muted}[1]{\textcolor{neutral}{#1}}
\newcommand{\key}[1]{\textcolor{primary-gold}{\textbf{#1}}}
\newcommand{\good}[1]{\textcolor{positive}{#1}}
\newcommand{\bad}[1]{\textcolor{negative}{#1}}

% Standout frame macro: full-bleed dark background for section transitions.
% Beamer-only (uses \begin{frame}/\setbeamercolor/\usebeamerfont) -- do not
% call from an article-class file (e.g. Notes/<CODE>/*-notes.tex). \muted,
% \key, \good, \bad, and \coursecode above are plain xcolor/gdef and are
% safe to use from either class; verified when Notes/article-class support
% was added.
% Expand: \transitionslide{Your transition title}
\newcommand{\transitionslide}[1]{%
  \begingroup
  \setbeamercolor{background canvas}{bg=primary-blue}
  \begin{frame}[plain]
    \centering
    \vfill
    {\usebeamerfont{title}\color{white}\Large #1\par}
    \vfill
  \end{frame}
  \endgroup
}

% Section divider frame (Beamer-only), an alternative to \transitionslide with a
% darker two-line style: near-black (jet) background, a small white label line
% above a larger gold title, and a thin gold rule along the bottom edge.
% Expand: \sectiondivider{Section 2}{Pointers}
\newcommand{\sectiondivider}[2]{%
  \begingroup
  \setbeamercolor{background canvas}{bg=jet}
  \begin{frame}[plain]
    \vspace*{3.0cm}
    {\color{white}\ttfamily\large #1\par}
    \vspace{0.5cm}
    {\color{primary-gold}\LARGE #2\par}
    \begin{tikzpicture}[remember picture, overlay]
      \draw[primary-gold, line width=2.5pt]
        ($(current page.south west)+(0,0.1)$) -- ($(current page.south east)+(0,0.1)$);
    \end{tikzpicture}
  \end{frame}
  \endgroup
}

% -----------------------------------------------------------------------------
% Channel watermark — "Concepts That Click" (YouTube).
%
% Article-class files (CompetitiveExam Q/A sets, Notes, Assignments) get a
% light diagonal draft watermark on every page plus a centered footer naming
% the channel. Beamer decks get a subtle TikZ background watermark instead
% (the footline already carries the course code). Centralized here so every
% MCQ PDF is branded without per-file edits.
% -----------------------------------------------------------------------------
\makeatletter
\@ifclassloaded{article}{%
  \usepackage{draftwatermark}
  \SetWatermarkText{Concepts That Click}
  \SetWatermarkScale{1}
  \SetWatermarkFontSize{2cm}
  \SetWatermarkLightness{0.86}
  \SetWatermarkAngle{45}
  \usepackage{fancyhdr}
  \pagestyle{fancy}
  \fancyhf{}
  \fancyfoot[C]{\footnotesize\textcolor{neutral}{Concepts That Click~|~YouTube: \href{https://www.youtube.com/@concepts.thatclick}{@concepts.thatclick}}}
  \renewcommand{\headrulewidth}{0pt}
  \renewcommand{\footrulewidth}{0pt}
  \fancypagestyle{plain}{%
    \fancyhf{}%
    \fancyfoot[C]{\footnotesize\textcolor{neutral}{Concepts That Click~|~YouTube: \href{https://www.youtube.com/@concepts.thatclick}{@concepts.thatclick}}}%
    \renewcommand{\headrulewidth}{0pt}%
    \renewcommand{\footrulewidth}{0pt}%
  }
}{}
\@ifclassloaded{beamer}{%
  \setbeamertemplate{background}{%
    \begin{tikzpicture}[remember picture, overlay]
      \node[rotate=45, opacity=0.055, align=center]
        at (current page.center)
        {\Huge\textcolor{neutral}{Concepts That Click}};
    \end{tikzpicture}%
  }
}{}
\makeatother 
\coursecode{CS301}
\renewcommand{\thesection}{5.\arabic{section}}

\title{Competitive Exam Practice 5: Queues, Circular Buffers, Deques, and Priority Queues\\\large GATE (Computer Science) Pattern}
\author{Auqib Hamid Lone\\[0.3em]
  \emph{Department of Computer Science \& Engineering}, \emph{GCET Ganderbal}}
\date{\today}

\begin{document}
\maketitle

\noindent Each question carries 1--2 marks. MCQ = single correct option; NAT = Numerical Answer Type.

\begin{enumerate}
\item[\textbf{Q1}] \textit{[GATE CSE 1996 Q1.12, verified --- GATEOverflow V2 p.240]} \textbf{[MCQ, 1 mark]} Consider the following statements:
\begin{itemize}
  \item[(i)] First-in-first out types of computations are efficiently supported by STACKS.
  \item[(ii)] Implementing LISTS on linked lists is more efficient than implementing LISTS on an array for almost all the basic LIST operations.
  \item[(iii)] Implementing QUEUES on a circular array is more efficient than implementing QUEUES on a linear array with two indices.
  \item[(iv)] Last-in-first-out type of computations are efficiently supported by QUEUES.
\end{itemize}
Which of the following is correct?
\begin{enumerate}
  \item[(A)] (ii) and (iii) are true
  \item[(B)] (i) and (ii) are true
  \item[(C)] (iii) and (iv) are true
  \item[(D)] (ii) and (iv) are true
\end{enumerate}

\item[\textbf{Q2}] \textit{[GATE CSE 2012 Q35, verified --- GATEOverflow V2 p.241]} \textbf{[MCQ, 2 marks]} Suppose a circular queue of capacity $(n-1)$ elements is implemented with an array of $n$ elements. Assume that the insertion and deletion operations are carried out using \texttt{REAR} and \texttt{FRONT} as array index variables, respectively. Initially, \texttt{REAR} = \texttt{FRONT} = 0. The conditions to detect queue full and queue empty are:
\begin{enumerate}
  \item[(A)] \texttt{full}: (\texttt{REAR}+1) mod $n$ == \texttt{FRONT} \quad \texttt{empty}: \texttt{REAR} == \texttt{FRONT}
  \item[(B)] \texttt{full}: (\texttt{REAR}+1) mod $n$ == \texttt{FRONT} \quad \texttt{empty}: (\texttt{FRONT}+1) mod $n$ == \texttt{REAR}
  \item[(C)] \texttt{full}: \texttt{REAR} == \texttt{FRONT} \quad \texttt{empty}: (\texttt{REAR}+1) mod $n$ == \texttt{FRONT}
  \item[(D)] \texttt{full}: (\texttt{FRONT}+1) mod $n$ == \texttt{REAR} \quad \texttt{empty}: \texttt{REAR} == \texttt{FRONT}
\end{enumerate}

\item[\textbf{Q3}] \textit{[GATE CSE 2016 Set 1 Q10, verified --- GATEOverflow V2 p.242]} \textbf{[MCQ, 1 mark]} A queue is implemented using an array such that ENQUEUE and DEQUEUE operations are performed efficiently. Which one of the following statements is CORRECT ($n$ refers to the number of items in the queue)?
\begin{enumerate}
  \item[(A)] Both operations can be performed in $O(1)$ time
  \item[(B)] At most one operation can be performed in $O(1)$ time but the worst case time for the other operation will be $\Omega(n)$
  \item[(C)] The worst case time complexity for both operations will be $\Omega(n)$
  \item[(D)] Worst case time complexity for both operations will be $\Omega(\log n)$
\end{enumerate}

\item[\textbf{Q4}] \textit{[GATE CSE 1997 Q4.7, verified --- GATEOverflow V2 p.239]} \textbf{[MCQ, 1 mark]} A priority queue \texttt{Q} is used to implement a stack that stores characters. \texttt{PUSH(C)} is implemented as \texttt{INSERT(Q, C, K)} where \texttt{K} is an appropriate integer key chosen by the implementation. \texttt{POP} is implemented as \texttt{DELETEMIN(Q)}. For a sequence of operations, the keys chosen are in
\begin{enumerate}
  \item[(A)] non-increasing order
  \item[(B)] non-decreasing order
  \item[(C)] strictly increasing order
  \item[(D)] strictly decreasing order
\end{enumerate}

\item[\textbf{Q5}] \textit{[GATE IT 2007 Q30, verified --- GATEOverflow V2 pp.243--244]} \textbf{[MCQ, 2 marks]} Suppose you are given an implementation of a queue of integers. The operations that can be performed on the queue are:
\begin{itemize}
  \item[(i)] \texttt{isEmpty(Q)} --- returns true if the queue is empty, false otherwise.
  \item[(ii)] \texttt{delete(Q)} --- deletes the element at the front of the queue and returns its value.
  \item[(iii)] \texttt{insert(Q, i)} --- inserts the integer \texttt{i} at the rear of the queue.
\end{itemize}
Consider the following function:
\begin{verbatim}
void f (queue Q) {
    int i;
    if (!isEmpty(Q)) {
        i = delete(Q);
        f(Q);
        insert(Q, i);
    }
}
\end{verbatim}
What operation is performed by the above function \texttt{f}?
\begin{enumerate}
  \item[(A)] Leaves the queue \texttt{Q} unchanged
  \item[(B)] Reverses the order of the elements in the queue \texttt{Q}
  \item[(C)] Deletes the element at the front of the queue \texttt{Q} and inserts it at the rear keeping the other elements in the same order
  \item[(D)] Empties the queue \texttt{Q}
\end{enumerate}

\item[\textbf{Q6}] \textit{[Original, GATE-pattern]} \textbf{[MCQ, 2 marks]} A circular queue of capacity $6$ uses the \texttt{count} version (\texttt{count} $= 0$ means empty, \texttt{count} $= 6$ means full), with indices advancing modulo $6$. At some moment \texttt{front} $= 4$, \texttt{rear} $= 2$, \texttt{count} $= 5$. One \texttt{dequeue} followed by one \texttt{enqueue} is performed. Afterwards, (\texttt{front}, \texttt{rear}, \texttt{count}) equals:
\begin{enumerate}
  \item[(A)] $(5, 3, 5)$
  \item[(B)] $(5, 3, 6)$
  \item[(C)] $(4, 3, 5)$
  \item[(D)] $(5, 2, 5)$
\end{enumerate}

\item[\textbf{Q7}] \textit{[Original, GATE-pattern]} \textbf{[MCQ, 1 mark]} A double-ended queue supports insertions at \emph{both} ends but deletions from the \emph{front} only. Which restricted variant is this, and which operation is forbidden?
\begin{enumerate}
  \item[(A)] Input-restricted deque; \texttt{enqueueFront} is forbidden
  \item[(B)] Output-restricted deque; \texttt{dequeueRear} is forbidden
  \item[(C)] Input-restricted deque; \texttt{dequeueRear} is forbidden
  \item[(D)] Output-restricted deque; \texttt{enqueueFront} is forbidden
\end{enumerate}

\item[\textbf{Q8}] \textit{[Original, GATE-pattern]} \textbf{[MCQ, 2 marks]} A priority queue holds up to $n = 10{,}000$ items in an array (no heap allowed). The workload is $10{,}000$ inserts followed by $5$ \texttt{deleteMin} operations. Which design gives the smaller worst-case total time?
\begin{enumerate}
  \item[(A)] Unordered array: $10{,}000$ $O(1)$ inserts plus $5$ $O(n)$ scans
  \item[(B)] Ordered array: $10{,}000$ $O(n)$ inserts plus $5$ $O(1)$ removals
  \item[(C)] Both designs cost the same asymptotically on this workload
  \item[(D)] Ordered array, because \texttt{deleteMin} dominates the insert cost
\end{enumerate}

\end{enumerate}

\end{document}
